\documentclass[../../main.tex]{subfiles}
\begin{document}

\begin{center}
  {\fontsize{48pt}{58pt}\selectfont \textbf{1995}\normalsize\textbf{年度}}
\end{center}

\begin{table}[htb]
  \centering
  \caption{出題テーマと難易度}
  \begin{tabular}{|c|p{9cm}|c|}
    \hline
    \textbf{問題番号} & \textbf{テーマ} & \textbf{難易度} \\
    \hline
    第1問 & 不等式の最良定数 & \\
    \hline
    第2問 & 積分で定められた関数の凸性 & \\
    \hline
    第3問 & 長方形の敷き詰めと漸化式 & \\
    \hline
    第4問 & 約数を用いた式の最小値 & \\
    \hline
    第5問 & サイコロによる平面上の移動と確率 & \\
    \hline
    第6問 & 双曲線の接線と三角形の面積 & \\
    \hline
  \end{tabular}
\end{table}

\begin{table}[htb]
  \centering
  \caption{解答}
  \renewcommand{\arraystretch}{1.4}
  \begin{tabular}{|c|c|p{9cm}|}
    \hline
    \textbf{問題番号} & \textbf{小問} & \textbf{答え} \\
    \hline
    第1問 & -- & $k=\dfrac{\sqrt{6}}{2}$ \\
    \hline
    第2問 & -- & $g(x+y)+g(x-y)\ge 2g(x)$（証明問題） \\
    \hline
    第3問 & (1) & $A_n=A_{n-1}+2A_{n-2}$ \\
    \cline{2-3}
    & (2) & $A_n=\dfrac{2^{n+1}+(-1)^n}{3}$ \\
    \hline
    第4問 & (1) & $k$が偶数なら$2\sqrt{N}$，奇数なら$\dfrac{3\sqrt{2N}}{2}$ \\
    \cline{2-3}
    & (2) & $142$ \\
    \hline
    第5問 & (1) & 途中で$x$軸を離れない（証明問題） \\
    \cline{2-3}
    & (2) & $\dfrac{1}{4}\left\{1-2\left(\dfrac{2}{3}\right)^n+\left(\dfrac{1}{3}\right)^n\right\}$ \\
    \hline
    第6問 & (1) & $S=ab$ \\
    \cline{2-3}
    & (2) & 最小値は$7\cdot 5^{-1/3}+5^{2/3}$ \\
    \hline
  \end{tabular}
\end{table}

\paragraph{試験形態}
6問．試験時間は要確認．

\end{document}
